% TOP_PROBLEM: {"id": 3000036, "title": "Finding kernels in special digraphs", "category_id": 2, "category": {"id": 2, "name": "combinatorics", "display_name": "Combinatorics", "description": "Counting problems, graph theory, discrete structures.", "slug": "combinatorics", "order_index": 2, "created_at": "2026-07-31T15:26:25.671Z"}, "status": "open", "problem_number": "AMR-029-0036", "set_id": 13}
% FIRST_POSTED: 2026-10-02
% ENTRY_KIND: statement_only
% UnsolvedMath ID 3000036; source catalogue code AMR-029-0036
% !TeX program = lualatex
% Definition and sources only; this file is not an LLM solution attempt.
\documentclass[11pt,a4paper]{article}
\usepackage[margin=25mm]{geometry}
\usepackage{fontspec}
\setmainfont{lmroman10-regular.otf}[BoldFont=lmroman10-bold.otf,
 ItalicFont=lmroman10-italic.otf,BoldItalicFont=lmroman10-bolditalic.otf]
\usepackage{amsmath,amssymb,mathtools}
\usepackage{unicode-math}
\setmathfont{latinmodern-math.otf}
\AtBeginDocument{\let\setminus\smallsetminus}
\usepackage{xurl,hyperref}
\hypersetup{colorlinks=true,urlcolor=blue}
\setcounter{secnumdepth}{0}
\setlength{\parindent}{0pt}
\setlength{\parskip}{5pt}
\setlength{\emergencystretch}{3em}
\begin{document}
\section*{Finding kernels in special digraphs}\label{3000036}
{\small UnsolvedMath ID 3000036 · AMR-029-0036 · Combinatorics · AMR Open Problem Lists}
\subsection{Definitions and mathematical statement}
Let $D=(V,A)$ be a finite loopless digraph. Opposite arcs may both be
present. A kernel is a set $K\subseteq V$ which is independent and
absorbing: no arc has both endpoints in $K$, and
\[
             \text{for every }v\notin K\text{ there is }u\in K
                         \text{ with }vu\in A.
\]
Thus absorption uses an arc from an outside vertex into the set, not
the reverse direction.

Identify structurally defined classes $\mathcal D$ for which kernel
existence can be decided, and a kernel constructed when one exists,
in time polynomial in $|V|+|A|$. The algorithm's exponent is fixed
for the class. State whether membership is promised or efficiently
recognizable and whether an additional representation is required.
The question does not require polynomial solvability for arbitrary
digraphs, where the decision problem is NP-complete.

A specific source-highlighted target is a super-orientation of a perfect
graph: replace each undirected edge by one or both directions. A graph
is perfect when every induced subgraph has chromatic number equal to
its largest clique size. Assume that within every clique, the one-way
arcs form an acyclic digraph. Under this hypothesis existence is known;
the remaining task is efficient construction for the whole class.
Results for narrower classes should specify their extra restrictions.
A kernel of one digraph is not a certificate that every induced
subdigraph has a kernel, which is the separate kernel-perfectness question.
\subsection{Short English statement}
Find structural classes of digraphs in which a kernel can be found or its absence certified in polynomial time, including the stated perfect-graph class.
\subsection{Sources}
\begin{itemize}
\item[S1] ULAM.AI and the UnsolvedMath Contributors, supplied problems.json, record 3000036 (AMR-029-0036), AMR Open Problem Lists. Catalogue identity and source transcription; CC BY 4.0.
\url{https://huggingface.co/datasets/ulamai/UnsolvedMath}.
\item[S2] Egres Open - List of open problems, Open-problem page: Finding kernels in special digraphs. Original source indexed by AMR-029-0036.
\url{https://oldlemon.cs.elte.hu/egres/open/List_of_open_problems}.
\item[S3] EGRES Open, Finding kernels in special digraphs. Individual source statement and remarks.
\url{https://lemon.cs.elte.hu/egres/open/Finding_kernels_in_special_digraphs}.
\end{itemize}
\end{document}
